% Part: second-order-logic
% Chapter: sol-and-sets
% Section: power-of-continuum

\documentclass[../../../include/open-logic-section]{subfiles}

\begin{document}

\olfileid{sol}{set}{pow}

\olsection{అవిచ్ఛిన్న సమితి పరిమాణం}

\begin{explain}
ద్వితీయ-స్థాయి తర్కంలో \tetoken{వ్యక్తి క్షేత్రపు}{!!{domain}} ఉపసమితులపై
పరిమాణీకరించవచ్చు; కానీ \tetoken{వ్యక్తి క్షేత్రపు}{!!{domain}} ఉపసమితుల సమితులపై
పరిమాణీకరించలేం. దీన్ని నేరుగా చేయడానికి \emph{తృతీయ-స్థాయి} తర్కం
అవసరం. ఉదాహరణకు, ఒక సమితి ఘాత సమితి నుంచి ఆ సమితికే
\tetoken{ఒకటి-ఒకటి}{!!{injective}} ప్రమేయం ఏదీ లేదనే కాంటర్ సిద్ధాంతాన్ని
చెప్పాలనుకుంటే, “ప్రతి సమితి~$X$కు, ప్రతి సమితి~$P$కు, $P$ అనేది
$X$ ఘాత సమితి అయితే, $\cardle{P}{X}$ కాదు” అని సూత్రీకరించడానికి
ప్రయత్నించవచ్చు. $P$ అనేది~$X$ ఘాత సమితి అని చెప్పడానికి,
$P$లోని \tetoken{మూలకాలు}{!!{element}s} సరిగ్గా~$X$ యొక్క ఉపసమితులన్నీ అని
అధికారికంగా చెప్పాలి; అంటే $\lforall[Y][(P(Y)
  \liff Y \subseteq X)]$ వంటిది కావాలి. సమస్య $P(Y)$లో ఉంది: అది
ద్వితీయ-స్థాయి తర్కపు \tetoken{ఒక సూత్రం}{!!a{formula}} కాదు; ఎందుకంటే~$P$ వంటి
ఏకస్థాన సంబంధ చరాలకు పదాలు మాత్రమే ప్రవేశాలుగా ఉండగలవు.

అయితే, వ్యక్తి క్షేత్రం తగినంత పెద్దదైతే, సమితుల సమితులపై
పరిమాణీకరణను \emph{అనుకరించవచ్చు}. ద్విస్థాన సంబంధం~$R$,
\tetoken{వ్యక్తి క్షేత్రపు}{!!{domain}} \tetoken{మూలకాలను}{!!{element}s}
\tetoken{వ్యక్తి క్షేత్రపు}{!!{domain}} \tetoken{మూలకాలతో}{!!{element}s} సంబంధపరుస్తుందనే
వాస్తవాన్ని ఉపయోగించుకోవడమే ఆలోచన. అటువంటి~$R$ ఇచ్చినప్పుడు, ఏదో~$x$
దేనితోనైతే $R$-సంబంధం కలిగి ఉందో ఆ \tetoken{మూలకాలన్నిటినీ}{!!{element}s}
సేకరించవచ్చు: $\Setabs{y \in
  \Domain{M}}{R(x,y)}$ అనేది~$x$ “సంకేతీకరించే” సమితి. ఎదురు దిశలో,
$Z \subseteq \Pow{\Domain{M}}$ అనేది~$\Domain{M}$ ఉపసమితుల ఏదైనా
సమూహమై,~$Z$లోని సమితులు ఎన్ని ఉన్నాయో కనీసం అన్ని
\tetoken{మూలకాలు}{!!{element}s}~$\Domain{M}$లో ఉంటే, ప్రతి $Y \in Z$ను
ఏదో~$x$తో~$R$ ద్వారా సంకేతీకరించే సంబంధం
$R \subseteq \Domain{M}^2$ కూడా ఉంటుంది.
\end{explain}

\begin{defn}
$R \subseteq \Domain{M}^2$ అయితే, \emph{$x$, $R$ ద్వారా
సంకేతీకరించేది} $\Setabs{y
  \in \Domain{M}}{R(x, y)}$.
\end{defn}

\tetoken{ఒక మూలకం}{!!a{element}}~$x \in \Domain{M}$, $R$ ద్వారా
$Z \subseteq \Domain{M}$ను సంకేతీకరిస్తే, $Y \subseteq \Domain{M}$
అనే సమితి ఒక సమితుల సమితిని, అంటే~$Y$లోని
\tetoken{మూలకాలు}{!!{element}s} సంకేతీకరించే సమితులను సంకేతీకరిస్తుంది.
కాబట్టి ఒక సమితి~$Y$, $R$ ద్వారా $\Pow{X}$ను సంకేతీకరించగలదు. ప్రతి
$Z \subseteq X$కు, $R$ ద్వారా~$Z$ను సంకేతీకరించే ఏదో $x \in Y$
ఉండి, ప్రతి $x \in Y$ కూడా~$R$ ద్వారా ఏదో $Z \subseteq X$ను
సంకేతీకరిస్తేనే అది అలా చేస్తుంది.

\begin{prop}
\tetoken{సూత్రం}{!!{formula}}
  \[
  \fn{Codes}(x, R, Z) \ident \lforall[y][(Z(y) \liff
    R(x, y))]
  \]
$s(x)$ అనేది $s(R)$ ద్వారా $s(Z)$ను సంకేతీకరిస్తుందని వ్యక్తం
చేస్తుంది. \tetoken{సూత్రం}{!!{formula}}
\begin{multline*}
  \fn{Pow}(Y, R, X) \ident \\
  \lforall[Z][(Z \subseteq X \lif \lexists[x][(Y(x) \land
    \fn{Codes}(x, R, Z))])] \land {} \\ \lforall[x][(Y(x) \lif
  \lforall[Z][(\fn{Codes}(x, R, Z) \lif Z \subseteq X)])]
\end{multline*}
$s(Y)$ అనేది $s(R)$ ద్వారా~$s(X)$ ఘాత సమితిని సంకేతీకరిస్తుందని,
అంటే $s(Y)$లోని మూలకాలు $s(R)$ ద్వారా $s(X)$ ఉపసమితులను సరిగ్గా
సంకేతీకరిస్తాయని వ్యక్తం చేస్తుంది.
\sourcecorrection{OLTESOLSET-003}{ఆంగ్ల మూలంలోని $\fn{Pow}(Y,R,X)$ సూత్రం రెండవ సార్వత్రిక పరిమాణకంలో $Y(x)$తో మొదలైన నియమితానికి ముగింపు కుండలీకరణను విడిచిపెట్టింది. ఆ నియమితాన్ని దాని పరిమాణకపు ముగింపు చతురస్ర కుండలీకరణకు ముందు మూసాం.}
\end{prop}

\begin{explain}
ఈ ఉపాయంతో, ఉపసమితులపైనే పరిమాణీకరించకుండా వాటి సంకేతాలపై
పరిమాణీకరించి ఘాత సమితి గురించిన ప్రకటనలను వ్యక్తం చేయవచ్చు.
ఉదాహరణకు,~$X$ ఘాత సమితిని సంకేతీకరించే ఏ సంబంధపు వ్యక్తి క్షేత్రం
నుంచీ~$X$కే \tetoken{ఒకటి-ఒకటి}{!!{injective}} ప్రమేయం లేదని చెప్పడం ద్వారా
ఇప్పుడు కాంటర్ సిద్ధాంతాన్ని వ్యక్తం చేయవచ్చు.
\end{explain}

\begin{prop}
ఈ వాక్యం
\begin{multline*}
  \lforall[X][\lforall[Y][\lforall[R][(\fn{Pow}(Y, R, X) \lif \\
      \lnot \lexists[u][(\lforall[x][\lforall[y][(\eq[u(x)][u(y)] \lif
            \eq[x][y])]] \land {}\\
        \lforall[x][(Y(x) \lif X(u(x)))])])]]]
\end{multline*}
చెల్లుబాటవుతుంది.
\end{prop}

\begin{explain}
\tetoken{అనంతంగా లెక్కించదగిన}{!!a{denumerable}} సమితి ఘాత సమితి
\tetoken{లెక్కించలేనిది}{!!{nonenumerable}}; కాబట్టి దాని కార్డినాలిటీ ఏ
\tetoken{అనంతంగా లెక్కించదగిన}{!!{denumerable}} సమితి కార్డినాలిటీ కంటే కూడా
పెద్దది (అది $\aleph_0$). $\Pow{\Nat}$ పరిమాణాన్ని “అవిచ్ఛిన్న
సమితి పరిమాణం” అంటారు; ఎందుకంటే అది వాస్తవ సంఖ్యారేఖపైని బిందువుల
సమితి~$\Real$ పరిమాణంతో సమానం. ఒక
\tetoken{అనంతంగా లెక్కించదగిన}{!!a{denumerable}} సమితి ఘాత సమితిని
సంకేతీకరించేంత పెద్దదిగా \tetoken{వ్యక్తి క్షేత్రం}{!!{domain}} ఉంటే, పరిమాణం
$\aleph_0$ గల సమితి~$X$ ఘాత సమితిని సంకేతీకరించే ఏ సమితి~$Y$తోనైనా
సమసంఖ్యాకంగా ఉందని చెప్పడం ద్వారా ఒక సమితి అవిచ్ఛిన్న సమితి పరిమాణం
గలదని వ్యక్తం చేయవచ్చు. (వ్యక్తి క్షేత్రం తగినంత పెద్దది కాకపోతే,
అంటే~$\Real$తో సమసంఖ్యాకమైన ఉపసమితి దానిలో లేకపోతే,
$\Pow{X}$ను సంకేతీకరించే సంబంధం కూడా ఉండదు.)
\end{explain}

\begin{prop}
$\cardle{\Real}{\Domain{M}}$ అయితే, \tetoken{సూత్రం}{!!{formula}}
\begin{multline*}
\fn{Cont}(Y) \ident
\lexists[X][\lexists[R][((\fn{Aleph}_0(X) \land
      \fn{Pow}(Y, R, X)) \land \\
      \lforall[x][\lforall[y][((Y(x) \land {}
      Y(y) \land \lforall[z][R(x,z) \liff R(y,z)]) \lif \eq[x][y])]])]]
\end{multline*}
$\cardeq{s(Y)}{\Real}$ అని వ్యక్తం చేస్తుంది.
\end{prop}

\begin{proof}
  $\fn{Pow}(Y, R, X)$ అనేది $s(Y)$, $s(R)$ ద్వారా~$s(X)$ ఘాత
  సమితిని సంకేతీకరిస్తుందని వ్యక్తం చేస్తుంది; $\fn{Aleph}_0(X)$
  ప్రకారం అది లెక్కించదగినది. కాబట్టి $s(Y)$ కనీసం అవిచ్ఛిన్న
  సమితి పరిమాణమంత పెద్దది; అయితే ఇంకా పెద్దదై ఉండవచ్చు ($s(Y)$లోని
  అనేక మూలకాలు~$X$ యొక్క ఒకే ఉపసమితిని సంకేతీకరిస్తే). చివరి
  సంయోజకం దీన్ని తొలగిస్తుంది; అది $s(R)$ ద్వారా $s(Y)$ మూలకాలకు,
  $s(X)$ ఉపసమితులకు మధ్య ఉన్న అనుబంధం
  \tetoken{ఒకటి-ఒకటిగా}{!!{injective}} ఉండాలని కోరుతుంది.
\sourcecorrection{OLTESOLSET-004}{చివరి సంయోజకం గురించి ఆంగ్ల మూల వివరణ $s(Y)$ మూలకాలను $s(Z)$ ఉపసమితులతో అనుబంధిస్తుందని చెప్పింది; కానీ $Z$ ఇక్కడ బద్ధ చరం మాత్రమే, ఘాత సమితికి ఆధార సమితి $X$. అందువల్ల నియంత్రక సూత్రానికి అనుగుణంగా $s(X)$ ఉపసమితులు అని రాశాం.}
\end{proof}

\begin{prop}
$\cardeq{\Domain{M}}{\Real}$ అయితేనే
\begin{multline*}
  \Sat{M}{\lexists[X][\lexists[Y][\lexists[R][(\fn{Aleph}_0(X) \land \fn{Pow}(Y, R, X) \land \\
          \lexists[u][(\lforall[x][\lforall[y][(\eq[u(x)][u(y)] \lif \eq[x][y])]] \land {} \\
          \lforall[x][Y(u(x))] \land {} \\\lforall[y][(Y(y) \lif \lexists[x][\eq[y][u(x)]])])])]]]}.
        \end{multline*}
\sourcecorrection{OLTESOLSET-005}{వ్యక్తి క్షేత్రం, అవిచ్ఛిన్న సమితి సమసంఖ్యాకాలని చెప్పాల్సిన ఆంగ్ల మూల సూత్రం $u$ను ఒకటి-ఒకటిగా, $Y$పైకి సంగ్రస్తంగా కోరింది; కానీ $u$ విలువలన్నీ $Y$లో ఉండాలని కోరలేదు. గుర్తింపు ప్రమేయమే ఆ బలహీన నిబంధనలను ఏ $Y$కైనా నెరవేర్చగలదు. $u$ పరిధి $Y$లోనే ఉండే నిబంధనను చేర్చి, $u$ను వ్యక్తి క్షేత్రం నుంచి $Y$పైకి ద్విగుణ ప్రమేయంగా చేశాం.}
  \end{prop}

\begin{explain}
అవిచ్ఛిన్న సమితి పరిమాణం మొదటి
\tetoken{లెక్కించలేని}{!!{nonenumerable}} కార్డినాలిటీ, అంటే
$\Pow{\Nat}$ పరిమాణం~$\aleph_1$ అనే ప్రకటననే అవిచ్ఛిన్న సమితి
పరికల్పన అంటారు.
\end{explain}

\begin{prop}
అవిచ్ఛిన్న సమితి పరికల్పన సత్యమైతేనే \[\fn{CH} \ident
\lforall[X][(\fn{Aleph}_1(X) \liff \fn{Cont}(X))]\] చెల్లుబాటవుతుంది.
\end{prop}

$\lnot \fn{CH}$ చెల్లుబాటైతేనే అవిచ్ఛిన్న సమితి పరికల్పన అసత్యమని
చెప్పడం నిజం కాదని గమనించండి. \tetoken{ఒక లెక్కించదగిన}{!!a{enumerable}} వ్యక్తి
క్షేత్రంలో పరిమాణం~$\aleph_1$ గల ఉపసమితులూ ఉండవు, అవిచ్ఛిన్న సమితి
పరిమాణం గల ఉపసమితులూ ఉండవు; కాబట్టి
\tetoken{ఒక లెక్కించదగిన}{!!a{enumerable}} వ్యక్తి క్షేత్రంలో
$\fn{CH}$ ఎల్లప్పుడూ సత్యమే. అయితే అవిచ్ఛిన్న సమితి పరికల్పన
అసత్యమైతేనే చెల్లుబాటయ్యే వేరొక వాక్యాన్ని ఇవ్వవచ్చు:

\begin{prop}
అవిచ్ఛిన్న సమితి పరికల్పన అసత్యమైతేనే \[\fn{NCH} \ident
\lforall[X][(\fn{Cont}(X) \lif \lexists[Y][(Y \subseteq X \land \lnot
    \fn{Count}(Y) \land \lnot \cardeq{X}{Y})])]\] చెల్లుబాటవుతుంది.
\end{prop}

\end{document}
